% Part: applied-modal-logics
% Chapter: epistemic-logic
% Section: properties-accessibility

\documentclass[../../../include/open-logic-section]{subfiles}

\begin{document}

\olfileid{aml}{el}{acc}

\olsection{ప్రాప్యత సంబంధాలు, జ్ఞానసంబంధ సూత్రాలు}

సాధారణ మోడల్ తర్కాల్లో చట్ర అనురూపత గురించి మనకు
ఇప్పటికే తెలిసినదాన్ని బట్టి, జ్ఞానసంబంధ వ్యాఖ్యానంలో
లాక్షణిక \tetoken{సూత్రాలు}{!!{formula}s} ఎలా ఉంటాయో
చూడాలనుకోవచ్చు. జ్ఞానసంబంధ తర్కాలను సాధారణంగా
S5లో వ్యాఖ్యానిస్తారని చెప్పాం. కాబట్టి వివిధ
జ్ఞానసంబంధ సూత్రాలు ఎలా వ్యక్తమవుతాయో, వాటికి
వివిధ చట్ర షరతులు ఎలా అనురూపమవుతాయో చూద్దాం.

సాధారణ మోడల్ తర్కంలో వేర్వేరు మోడల్
\tetoken{సూత్రాలు}{!!{formula}s} ప్రాప్యత సంబంధాల
వేర్వేరు ధర్మాలను లక్షణీకరించాయని గుర్తుంచుకోండి.
ప్రత్యేక జ్ఞానసంబంధ సూత్రాలకు అనురూపమైన వాటిలో
కొన్నింటిని ఈ పట్టిక ఎంచి చూపుతుంది.

\begin{table}[t]
    \begin{tabular}{| p{.48\textwidth} || p{.48\textwidth} |}
      \hline
      {\emph{$R$ ఇలా ఉంటే \dots}} & {\emph{అప్పుడు~$\mModel{M}$లో ఇది సత్యం:}} \\
      \hline \hline

      & $\Knows (p \lif q) \lif (\Knows p \lif \Knows q)$
      \hfill \newline{(మూసుకుపోవడం)} \\
      \hline
      \emph{స్వావర్తనం}: $\forall w Rww$
      & $\Knows p \lif p$ \hfill \newline{(యథార్థసత్యత్వం)} \\
      \hline
      \emph{సంక్రామకం}: \newline
      $\forall u \forall v \forall w ((Ruv \land Rvw) \lif Ruw)$ &
      $\Knows p \lif \Knows \Knows p$ \hfill
           \newline{(సకారాత్మక స్వపరిశీలన)} \\
      \hline
      \emph{యూక్లిడియన్}: \newline
      $\forall w \forall u \forall v ((Rwu \land Rwv) \lif Ruv)$ &
      $\lnot \Knows p \lif \Knows \neg \Knows p$ \hfill
        \newline{(నకారాత్మక స్వపరిశీలన)} \\
      \hline
    \end{tabular}
    \caption{నాలుగు జ్ఞానసంబంధ సూత్రాలు.}
    \ollabel{tab:four}
  \end{table}

$T$ స్వీకృతానికి అనురూపమైన యథార్థసత్యత్వాన్ని
ఈ సూత్రాల్లో సాధారణంగా అతి తక్కువ వివాదాస్పదమైనదిగా
పరిగణిస్తారు. \tetoken{ఒక సూత్రం}{!!a{formula}}
తెలిసినదైతే అది సత్యమే కావాలనే వాదనను ఇది
సూచిస్తుంది. మూసుకుపోవడం, సకారాత్మక, నకారాత్మక
స్వపరిశీలనలు మరింత వివాదాస్పదమైనవి.

$K$ స్వీకృతానికి అనురూపమైన మూసుకుపోవడం, ఒక కర్త
జ్ఞానం అన్వయం కింద మూసుకుపోయి ఉంటుందనే భావనను
సూచిస్తుంది. కొన్ని సందర్భాల్లో ఇది సమంజసంగా
అనిపించవచ్చు. ఉదాహరణకు, నేను విక్టోరియాలో ఉంటే
వాంకూవర్ దీవిలో ఉంటాననే విషయం నాకు తెలిసి ఉండవచ్చు.
విచిత్రమైన సందేహవాద సందర్భాలను పక్కన పెడితే,
నేను విక్టోరియాలో ఉన్నానని నాకు తెలుసు; అందువల్ల
నేను వాంకూవర్ దీవిలో ఉన్నాననీ నాకు తెలుసని
భావించడం సహజం. ఈ సందర్భంలో నా జ్ఞానపు తార్కిక
మూసుకుపోవడం సహజంగానే తోచవచ్చు. మరోవైపు,
మన జ్ఞానపు పర్యవసానాలన్నింటినీ మనం ఎప్పుడూ
ఆలోచించి చూడం; కాబట్టి ఇతర సందర్భాల్లో ఇది
అంత సహజంగా అనిపించని ఫలితాలకు దారితీయవచ్చు.

కొన్నిసార్లు KK-సూత్రం అని పిలిచే సకారాత్మక
స్వపరిశీలనను ``నాకు ఏదైనా తెలిస్తే, అది నాకు
తెలుసని కూడా నాకు తెలుసు'' అని చెబుతారు.
ఇది 4 స్వీకృతానికి జ్ఞానసంబంధ అనురూపం.
అదే విధంగా నకారాత్మక స్వపరిశీలనను ``నాకు
ఏదైనా \emph{తెలియకపోతే}, అది నాకు తెలియదని
నాకు తెలుసు'' అని చెబుతారు; ఇది 5 స్వీకృతానికి
అనురూపం. ఈ రెండింటికీ సాధారణ ప్రతివుదాహరణలు
ఉన్నట్టు కనిపిస్తుంది: నాకు నిజంగా తెలిసిన
విషయం నాకు తెలుసో లేదో అనిశ్చితంగా ఉండవచ్చు.


\end{document}
